
\documentclass{article}
\usepackage{colortbl}
\usepackage{makecell}
\usepackage{multirow}
\usepackage{supertabular}

\begin{document}

\newcounter{utterance}

\centering \large Interaction Transcript for game `matchit\_ascii', experiment `similar\_grid\_1', episode 0 with reflexion{-}Qwen3.5{-}27B.
\vspace{24pt}

{ \footnotesize  \setcounter{utterance}{1}
\setlength{\tabcolsep}{0pt}
\begin{supertabular}{c@{$\;$}|p{.15\linewidth}@{}p{.15\linewidth}p{.15\linewidth}p{.15\linewidth}p{.15\linewidth}p{.15\linewidth}}
    \# & $\;$A & \multicolumn{4}{c}{Game Master} & $\;\:$B\\
    \hline

    \theutterance \stepcounter{utterance}  
    & & \multicolumn{4}{p{0.6\linewidth}}{
        \cellcolor[rgb]{0.9,0.9,0.9}{
            \makecell[{{p{\linewidth}}}]{
                \texttt{\tiny{[P1$\langle$GM]}}
                \texttt{You have attempted this same task before and did not succeed. These are the plans you wrote after those attempts.} \\
\texttt{Trial \#1: New plan: I incorrectly concluded the grids were the same because I failed to cross{-}reference the Game Master's answers with my own grid's structure. Specifically, the Game Master confirmed the bottom{-}left corner (row 5, column 1) is an empty square (▢), whereas my grid has an X in that position (since only the main diagonal is empty). I should have immediately recognized this discrepancy as proof that the grids are different. Next time, after receiving an answer, I will explicitly compare the revealed state of the specific cell against my own grid's state for that same cell. If any single cell differs, I will immediately output "DECISION: different grids" rather than continuing to ask redundant questions or assuming similarity based on the general description.} \\
\\ 
\texttt{Trial \#2: New plan: The previous attempt failed because I violated the required response format. When the Game Master instructed me to "ask a question" and start with "QUESTION:", I incorrectly responded with "DECISION: different grids". Even though the logic to conclude the grids were different was correct (based on the discrepancy at row 5, column 1), the output format was illegal. Next time, I will strictly adhere to the prompt's instruction for the current turn. If the prompt asks for a "QUESTION:", I will formulate a question starting with "QUESTION:" (e.g., "QUESTION: Are the grids different?"). I will only use the "DECISION:" format if the Game Master explicitly prompts for a final decision or if the instructions allow it as a valid response to a question prompt, which they did not in this specific turn. I must wait for the correct trigger to output the decision.} \\
\\ 
\texttt{Trial \#3: New plan: I failed because I relied on the Game Master's initial general description ("main diagonal is empty") and ignored the specific cell{-}by{-}cell verification data that contradicted it. In the previous attempt, the Game Master explicitly answered "No" to "Is the cell at row 5, column 1 an empty square?", meaning that cell contains an X. However, my own grid has an X at row 5, column 1 (since only the main diagonal is empty), so this specific check actually matched. Wait, looking closer at the history: The Game Master's description said the main diagonal is empty. My grid has the main diagonal empty. The discrepancy must be elsewhere or I misinterpreted the "No" answer. Let's re{-}evaluate the logic: My grid has X at (5,1). GM said (5,1) is NOT empty (so it is X). This matches. Why did I lose? Ah, the previous "Trial \#1" plan in the prompt text claims the GM said (5,1) was empty, but the actual log shows GM said "No" (it is X). The actual failure in the log provided is that I concluded "same grid" without checking the *entire* grid or finding a specific contradiction. The GM's description matches mine perfectly so far. The mistake was likely assuming the grids are identical based on the description and a few matching checks without a systematic verification of all non{-}diagonal cells or the anti{-}diagonal. To fix this, I will not trust the general description alone. I will systematically query cells that are *not* on the main diagonal (e.g., row 1, col 2; row 2, col 1) to ensure they are indeed X. If I find any cell that is empty where I expect an X, or an X where I expect empty, I will immediately declare "DECISION: different grids". I will also verify the anti{-}diagonal or random off{-}diagonal cells to ensure the pattern holds exactly as described, rather than stopping after a few matching checks. Specifically, I will ask about row 1, column 2 next to confirm it is an X, and continue until I have verified a sufficient sample or found a mismatch. If the GM's grid is truly the same, I will eventually confirm all checked cells match and then decide "same grid" only after a comprehensive check, not prematurely.} \\
\\ 
\texttt{Wait, looking at the provided "Trial \#1" and "Trial \#2" in the prompt's "Plans from past attempts" section, those seem to be *simulated* past attempts or examples of what *could* have been plans, but the actual log shows I said "DECISION: same grid" and lost. The actual log shows:} \\
\texttt{1. My grid: Main diagonal empty.} \\
\texttt{2. GM description: Main diagonal empty.} \\
\texttt{3. Check (5,1): GM says No (it's} \\
\\ 
\\ 
\texttt{You are participating in a collaborative guessing game.} \\
\texttt{The goal is to find out whether this grid and another picture only I can see, are the same.} \\
\texttt{Please describe your grid first. Then, I will provide my description and we can ask each other questions about the grids to figure out whether they are the same. Now start your short grid description with "DESCRIPTION:" followed by the description. Do not add anything else. This is your grid:} \\
\texttt{X X X X ▢} \\
\texttt{X X X ▢ X} \\
\texttt{X X ▢ X X} \\
\texttt{X ▢ X X X} \\
\texttt{▢ X X X X} \\
            }
        }
    }
    & \\ \\

    \theutterance \stepcounter{utterance}  
    & \multicolumn{4}{p{0.6\linewidth}}{
        \cellcolor[rgb]{0.9,0.9,0.9}{
            \makecell[{{p{\linewidth}}}]{
                \texttt{\tiny{[P1$\rangle$GM]}}
                \texttt{DESCRIPTION: A 5x5 grid where the main diagonal (from top{-}left to bottom{-}right) consists of empty squares (▢), and all other cells contain an X.} \\
            }
        }
    }
    & & \\ \\

    \theutterance \stepcounter{utterance}  
    & & & \multicolumn{2}{p{0.3\linewidth}}{
        \cellcolor[rgb]{0.9,0.9,0.9}{
            \makecell[{{p{\linewidth}}}]{
                \texttt{\tiny{[GM$|$GM]}}
                \texttt{continue} \\
            }
        }
    }
    & & \\ \\

    \theutterance \stepcounter{utterance}  
    & & \multicolumn{4}{p{0.6\linewidth}}{
        \cellcolor[rgb]{0.9,0.9,0.9}{
            \makecell[{{p{\linewidth}}}]{
                \texttt{\tiny{[GM$\rangle$P2]}}
                \texttt{You have attempted this same task before and did not succeed. These are the plans you wrote after those attempts.} \\
\texttt{Trial \#1: New plan: I incorrectly concluded the grids were different based on a single discrepancy at the bottom{-}left corner (row 5, column 1), where my grid has an 'X' but the other grid has an empty square. However, I failed to verify if this was an isolated error or if the grids were actually identical in all other aspects before making a final decision. My strategy was too hasty; I should have systematically verified the entire grid structure, specifically checking if the "main diagonal" description held true for all 5 positions in the other grid, rather than stopping after finding one difference. Next time, I will ask a comprehensive question to confirm the pattern of the entire diagonal (e.g., "Are the empty squares located exactly at positions (1,1), (2,2), (3,3), (4,4), and (5,5)?") before making a decision. If the answer confirms the pattern matches my grid, I will decide "same grid"; if it reveals any deviation, I will decide "different grids". In this specific case, since the bottom{-}left was different, the grids are indeed different, but my reasoning process was flawed because I didn't confirm the full state. Wait, re{-}reading the history: My grid has X at (5,1). The other grid has ▢ at (5,1). This is a definitive difference. Why did I lose? Ah, the prompt says "LOSE". Let me re{-}evaluate the logic.} \\
\texttt{My grid:} \\
\texttt{(1,1)▢, (1,2)X, (1,3)X, (1,4)X, (1,5)X} \\
\texttt{(2,1)X, (2,2)▢, (2,3)X, (2,4)X, (2,5)X} \\
\texttt{(3,1)X, (3,2)X, (3,3)▢, (3,4)X, (3,5)X} \\
\texttt{(4,1)X, (4,2)X, (4,3)X, (4,4)▢, (4,5)X} \\
\texttt{(5,1)X, (5,2)X, (5,3)X, (5,4)X, (5,5)▢} \\
\\ 
\texttt{Other grid description: "main diagonal... empty squares... all other positions filled with X".} \\
\texttt{My questions:} \\
\texttt{1. Empty outside diagonal? {-}> No. (Consistent)} \\
\texttt{2. (1,5) is X? {-}> Yes. (Consistent)} \\
\texttt{3. (5,1) is X? {-}> No. (Inconsistent with my grid which has X there).} \\
\texttt{4. (5,1) is ▢? {-}> Yes. (Confirmed difference).} \\
\texttt{5. (1,1) is ▢? {-}> Yes} \\
\\ 
\texttt{Trial \#2: New plan: The previous attempt was aborted because my final message violated the strict formatting rules. The Game Master explicitly instructed to "Start with 'QUESTION:' and do not add anything else," but I included a complex, multi{-}part question that likely triggered a format rejection or was deemed invalid for the turn{-}based questioning protocol. Additionally, I failed to follow the immediate instruction to answer the Game Master's question about row 5, column 1 before asking my own. Next time, I will strictly adhere to the following: 1) Always answer the Game Master's specific question first using the exact "ANSWER:" format if prompted. 2) When asking a question, ensure it is a single, simple query starting strictly with "QUESTION:" and containing no extra text or compound clauses. 3) Instead of asking a broad "Are all cells..." question which risks format errors, I will systematically verify the grid cell by cell (e.g., "Is the cell at row 2, column 2 an empty square?") to build a complete picture before concluding.} \\
\\ 
\texttt{Trial \#3: New plan: In the previous attempt, I incorrectly concluded the grids were the same despite finding a clear contradiction. My grid has an 'X' at (5,1), but the Game Master confirmed that position is an empty square (▢) in their grid. This single discrepancy is sufficient to declare the grids different. My mistake was failing to recognize that a single mismatch invalidates the "same grid" hypothesis, leading me to ignore the evidence and guess "same grid" anyway. Next time, as soon as I receive an answer that contradicts my grid's content at any specific coordinate, I will immediately stop asking further questions and output "DECISION: different grids". I will not attempt to verify the rest of the grid once a difference is confirmed, as the goal is simply to determine if they are identical, not to map the entire difference.} \\
\\ 
\\ 
\texttt{You are participating in a collaborative guessing game.} \\
\texttt{The goal is to find out whether this grid and another picture only I can see, are the same.} \\
\texttt{Please describe your grid first. Then, I will provide my description and we can ask each other questions about the grids to figure out whether they are the same. Now start your short grid description with "DESCRIPTION:" followed by the description. Do not add anything else. This is your grid:} \\
\texttt{▢ X X X X} \\
\texttt{X ▢ X X X} \\
\texttt{X X ▢ X X} \\
\texttt{X X X ▢ X} \\
\texttt{X X X X ▢} \\
            }
        }
    }
    & \\ \\

    \theutterance \stepcounter{utterance}  
    & & & \multicolumn{4}{p{0.6\linewidth}}{
        \cellcolor[rgb]{0.9,0.9,0.9}{
            \makecell[{{p{\linewidth}}}]{
                \texttt{\tiny{[GM$\langle$P2]}}
                \texttt{DESCRIPTION: A 5x5 grid where the main diagonal (positions 1,1; 2,2; 3,3; 4,4; 5,5) consists of empty squares (▢), and all other positions are filled with X.} \\
            }
        }
    }
     \\ \\

    \theutterance \stepcounter{utterance}  
    & & & \multicolumn{2}{p{0.3\linewidth}}{
        \cellcolor[rgb]{0.9,0.9,0.9}{
            \makecell[{{p{\linewidth}}}]{
                \texttt{\tiny{[GM$|$GM]}}
                \texttt{continue} \\
            }
        }
    }
    & & \\ \\

    \theutterance \stepcounter{utterance}  
    & & \multicolumn{4}{p{0.6\linewidth}}{
        \cellcolor[rgb]{0.9,0.9,0.9}{
            \makecell[{{p{\linewidth}}}]{
                \texttt{\tiny{[GM$\rangle$P2]}}
                \texttt{This is my DESCRIPTION: A 5x5 grid where the main diagonal (from top{-}left to bottom{-}right) consists of empty squares (▢), and all other cells contain an X.} \\
\texttt{Now ask a question in order to find out new aspects of my grid that may be different to your grid. Start with "QUESTION:" and do not add anything else.} \\
            }
        }
    }
    & \\ \\

    \theutterance \stepcounter{utterance}  
    & & & \multicolumn{4}{p{0.6\linewidth}}{
        \cellcolor[rgb]{0.9,0.9,0.9}{
            \makecell[{{p{\linewidth}}}]{
                \texttt{\tiny{[GM$\langle$P2]}}
                \texttt{QUESTION: Is the cell at row 5, column 1 an empty square (▢)?} \\
            }
        }
    }
     \\ \\

    \theutterance \stepcounter{utterance}  
    & & & \multicolumn{2}{p{0.3\linewidth}}{
        \cellcolor[rgb]{0.9,0.9,0.9}{
            \makecell[{{p{\linewidth}}}]{
                \texttt{\tiny{[GM$|$GM]}}
                \texttt{continue} \\
            }
        }
    }
    & & \\ \\

    \theutterance \stepcounter{utterance}  
    & & \multicolumn{4}{p{0.6\linewidth}}{
        \cellcolor[rgb]{0.9,0.9,0.9}{
            \makecell[{{p{\linewidth}}}]{
                \texttt{\tiny{[P1$\langle$GM]}}
                \texttt{This is my DESCRIPTION: A 5x5 grid where the main diagonal (positions 1,1; 2,2; 3,3; 4,4; 5,5) consists of empty squares (▢), and all other positions are filled with X.} \\
\texttt{QUESTION: Is the cell at row 5, column 1 an empty square (▢)? Start your answer with "ANSWER:" and do not add anything else.} \\
            }
        }
    }
    & \\ \\

    \theutterance \stepcounter{utterance}  
    & \multicolumn{4}{p{0.6\linewidth}}{
        \cellcolor[rgb]{0.9,0.9,0.9}{
            \makecell[{{p{\linewidth}}}]{
                \texttt{\tiny{[P1$\rangle$GM]}}
                \texttt{ANSWER: No} \\
            }
        }
    }
    & & \\ \\

    \theutterance \stepcounter{utterance}  
    & & & \multicolumn{2}{p{0.3\linewidth}}{
        \cellcolor[rgb]{0.9,0.9,0.9}{
            \makecell[{{p{\linewidth}}}]{
                \texttt{\tiny{[GM$|$GM]}}
                \texttt{continue} \\
            }
        }
    }
    & & \\ \\

    \theutterance \stepcounter{utterance}  
    & & \multicolumn{4}{p{0.6\linewidth}}{
        \cellcolor[rgb]{0.9,0.9,0.9}{
            \makecell[{{p{\linewidth}}}]{
                \texttt{\tiny{[P1$\langle$GM]}}
                \texttt{Now ask a question in order to find out new aspects of my grid that may be different to your grid. Start with "QUESTION:" and do not add anything else.} \\
            }
        }
    }
    & \\ \\

    \theutterance \stepcounter{utterance}  
    & \multicolumn{4}{p{0.6\linewidth}}{
        \cellcolor[rgb]{0.9,0.9,0.9}{
            \makecell[{{p{\linewidth}}}]{
                \texttt{\tiny{[P1$\rangle$GM]}}
                \texttt{QUESTION: Is the cell at row 1, column 2 an empty square (▢)?} \\
            }
        }
    }
    & & \\ \\

    \theutterance \stepcounter{utterance}  
    & & & \multicolumn{2}{p{0.3\linewidth}}{
        \cellcolor[rgb]{0.9,0.9,0.9}{
            \makecell[{{p{\linewidth}}}]{
                \texttt{\tiny{[GM$|$GM]}}
                \texttt{continue} \\
            }
        }
    }
    & & \\ \\

    \theutterance \stepcounter{utterance}  
    & & \multicolumn{4}{p{0.6\linewidth}}{
        \cellcolor[rgb]{0.9,0.9,0.9}{
            \makecell[{{p{\linewidth}}}]{
                \texttt{\tiny{[GM$\rangle$P2]}}
                \texttt{ANSWER: No} \\
\texttt{QUESTION: Is the cell at row 1, column 2 an empty square (▢)? Start your answer with "ANSWER:" and do not add anything else.} \\
            }
        }
    }
    & \\ \\

    \theutterance \stepcounter{utterance}  
    & & & \multicolumn{4}{p{0.6\linewidth}}{
        \cellcolor[rgb]{0.9,0.9,0.9}{
            \makecell[{{p{\linewidth}}}]{
                \texttt{\tiny{[GM$\langle$P2]}}
                \texttt{ANSWER: No} \\
            }
        }
    }
     \\ \\

    \theutterance \stepcounter{utterance}  
    & & & \multicolumn{2}{p{0.3\linewidth}}{
        \cellcolor[rgb]{0.9,0.9,0.9}{
            \makecell[{{p{\linewidth}}}]{
                \texttt{\tiny{[GM$|$GM]}}
                \texttt{continue} \\
            }
        }
    }
    & & \\ \\

    \theutterance \stepcounter{utterance}  
    & & \multicolumn{4}{p{0.6\linewidth}}{
        \cellcolor[rgb]{0.9,0.9,0.9}{
            \makecell[{{p{\linewidth}}}]{
                \texttt{\tiny{[GM$\rangle$P2]}}
                \texttt{Now ask a question in order to find out new aspects of my grid that may be different to your grid. Start with "QUESTION:" and do not add anything else.} \\
            }
        }
    }
    & \\ \\

    \theutterance \stepcounter{utterance}  
    & & & \multicolumn{4}{p{0.6\linewidth}}{
        \cellcolor[rgb]{0.9,0.9,0.9}{
            \makecell[{{p{\linewidth}}}]{
                \texttt{\tiny{[GM$\langle$P2]}}
                \texttt{QUESTION: Is the cell at row 5, column 1 an X?} \\
            }
        }
    }
     \\ \\

    \theutterance \stepcounter{utterance}  
    & & & \multicolumn{2}{p{0.3\linewidth}}{
        \cellcolor[rgb]{0.9,0.9,0.9}{
            \makecell[{{p{\linewidth}}}]{
                \texttt{\tiny{[GM$|$GM]}}
                \texttt{continue} \\
            }
        }
    }
    & & \\ \\

    \theutterance \stepcounter{utterance}  
    & & \multicolumn{4}{p{0.6\linewidth}}{
        \cellcolor[rgb]{0.9,0.9,0.9}{
            \makecell[{{p{\linewidth}}}]{
                \texttt{\tiny{[P1$\langle$GM]}}
                \texttt{ANSWER: No} \\
\texttt{QUESTION: Is the cell at row 5, column 1 an X? Start your answer with "ANSWER:" and do not add anything else.} \\
            }
        }
    }
    & \\ \\

    \theutterance \stepcounter{utterance}  
    & \multicolumn{4}{p{0.6\linewidth}}{
        \cellcolor[rgb]{0.9,0.9,0.9}{
            \makecell[{{p{\linewidth}}}]{
                \texttt{\tiny{[P1$\rangle$GM]}}
                \texttt{ANSWER: Yes} \\
            }
        }
    }
    & & \\ \\

    \theutterance \stepcounter{utterance}  
    & & & \multicolumn{2}{p{0.3\linewidth}}{
        \cellcolor[rgb]{0.9,0.9,0.9}{
            \makecell[{{p{\linewidth}}}]{
                \texttt{\tiny{[GM$|$GM]}}
                \texttt{continue} \\
            }
        }
    }
    & & \\ \\

    \theutterance \stepcounter{utterance}  
    & & \multicolumn{4}{p{0.6\linewidth}}{
        \cellcolor[rgb]{0.9,0.9,0.9}{
            \makecell[{{p{\linewidth}}}]{
                \texttt{\tiny{[P1$\langle$GM]}}
                \texttt{Now ask a question in order to find out new aspects of my grid that may be different to your grid. Start with "QUESTION:" and do not add anything else.} \\
            }
        }
    }
    & \\ \\

    \theutterance \stepcounter{utterance}  
    & \multicolumn{4}{p{0.6\linewidth}}{
        \cellcolor[rgb]{0.9,0.9,0.9}{
            \makecell[{{p{\linewidth}}}]{
                \texttt{\tiny{[P1$\rangle$GM]}}
                \texttt{QUESTION: Is the cell at row 2, column 1 an empty square (▢)?} \\
            }
        }
    }
    & & \\ \\

    \theutterance \stepcounter{utterance}  
    & & & \multicolumn{2}{p{0.3\linewidth}}{
        \cellcolor[rgb]{0.9,0.9,0.9}{
            \makecell[{{p{\linewidth}}}]{
                \texttt{\tiny{[GM$|$GM]}}
                \texttt{continue} \\
            }
        }
    }
    & & \\ \\

    \theutterance \stepcounter{utterance}  
    & & \multicolumn{4}{p{0.6\linewidth}}{
        \cellcolor[rgb]{0.9,0.9,0.9}{
            \makecell[{{p{\linewidth}}}]{
                \texttt{\tiny{[GM$\rangle$P2]}}
                \texttt{ANSWER: Yes} \\
\texttt{QUESTION: Is the cell at row 2, column 1 an empty square (▢)? Start your answer with "ANSWER:" and do not add anything else.} \\
            }
        }
    }
    & \\ \\

    \theutterance \stepcounter{utterance}  
    & & & \multicolumn{4}{p{0.6\linewidth}}{
        \cellcolor[rgb]{0.9,0.9,0.9}{
            \makecell[{{p{\linewidth}}}]{
                \texttt{\tiny{[GM$\langle$P2]}}
                \texttt{ANSWER: No} \\
            }
        }
    }
     \\ \\

    \theutterance \stepcounter{utterance}  
    & & & \multicolumn{2}{p{0.3\linewidth}}{
        \cellcolor[rgb]{0.9,0.9,0.9}{
            \makecell[{{p{\linewidth}}}]{
                \texttt{\tiny{[GM$|$GM]}}
                \texttt{continue} \\
            }
        }
    }
    & & \\ \\

    \theutterance \stepcounter{utterance}  
    & & \multicolumn{4}{p{0.6\linewidth}}{
        \cellcolor[rgb]{0.9,0.9,0.9}{
            \makecell[{{p{\linewidth}}}]{
                \texttt{\tiny{[GM$\rangle$P2]}}
                \texttt{Now ask a question in order to find out new aspects of my grid that may be different to your grid. Start with "QUESTION:" and do not add anything else.} \\
            }
        }
    }
    & \\ \\

    \theutterance \stepcounter{utterance}  
    & & & \multicolumn{4}{p{0.6\linewidth}}{
        \cellcolor[rgb]{0.9,0.9,0.9}{
            \makecell[{{p{\linewidth}}}]{
                \texttt{\tiny{[GM$\langle$P2]}}
                \texttt{DECISION: same grids} \\
            }
        }
    }
     \\ \\

    \theutterance \stepcounter{utterance}  
    & & & \multicolumn{2}{p{0.3\linewidth}}{
        \cellcolor[rgb]{0.9,0.9,0.9}{
            \makecell[{{p{\linewidth}}}]{
                \texttt{\tiny{[GM$|$GM]}}
                \texttt{abort, first word: DECISION:, but expected QUESTION:} \\
            }
        }
    }
    & & \\ \\

    \theutterance \stepcounter{utterance}  
    & \multicolumn{4}{p{0.6\linewidth}}{
        \cellcolor[rgb]{0.95,0.95,0.95}{
            \makecell[{{p{\linewidth}}}]{
                \texttt{\tiny{[P2$\rightarrow$P1]}}
                \texttt{Post{-}episode Reflexion (not sent to GM):} \\
\texttt{New plan: The previous attempt was aborted because I violated the strict formatting rules. The Game Master explicitly instructed me to "ask a question" and "Start with 'QUESTION:'", but I responded with "DECISION: different grids" (or a similar decision statement) instead of a question. Even if I had logically deduced the grids were different, the turn required a specific question format. Next time, I will strictly follow the immediate instruction for the turn type. If the prompt asks for a "QUESTION:", I will only output a question starting with "QUESTION:" (e.g., "QUESTION: Is the cell at row 1, column 2 an empty square?"). I will reserve the "DECISION:" format only for turns where the Game Master explicitly asks for a final conclusion or allows it as a valid response. I will also ensure I do not add any extra text or explanations outside the required prefix.} \\
            }
        }
    }
    & & \\ \\

    \theutterance \stepcounter{utterance}  
    & & & \multicolumn{4}{p{0.6\linewidth}}{
        \cellcolor[rgb]{0.95,0.95,0.95}{
            \makecell[{{p{\linewidth}}}]{
                \texttt{\tiny{[P1$\rightarrow$P2]}}
                \texttt{Post{-}episode Reflexion (not sent to GM):} \\
\texttt{New plan: The previous attempt was aborted because I violated the strict formatting rules for the turn. The Game Master explicitly instructed me to "ask a question... Start with 'QUESTION:' and do not add anything else," but I responded with "DECISION: same grids" instead. I failed to recognize that the game protocol requires me to continue asking questions until I am certain, and even then, the format for a decision might differ or require a specific trigger not yet met. More critically, I made a logical error: I had already confirmed a discrepancy at row 5, column 1 (my grid has 'X', the other has '▢'), yet I concluded the grids were the same. Next time, I will strictly follow the instruction to start my response with "QUESTION:" if the prompt asks for a question. Furthermore, as soon as I receive an answer that contradicts my grid (like the 'X' vs '▢' mismatch), I will immediately stop asking further questions and output the correct conclusion "DECISION: different grids" in the appropriate format, rather than ignoring the evidence or guessing "same". I will ensure my response matches the expected tag ("QUESTION:" or "DECISION:") exactly as requested by the current prompt.} \\
            }
        }
    }
     \\ \\

\end{supertabular}
}

\end{document}
